\documentclass[10pt,a4paper]{amsart}
\usepackage[margin=19mm]{geometry}
\usepackage{amsmath,amssymb,mathtools}
\usepackage[scr=rsfs,cal=euler]{mathalfa}
\usepackage{xfrac}
\usepackage[unicode,hidelinks,bookmarks=false]{hyperref}
\newcommand{\ie}{i.e.\ }
\newcommand{\cf}{cf.\ }
\newcommand{\resp}{respectively\ }
\newcommand{\Subsection}{\subsection}
\providecommand{\nobreakdash}{\nobreak\hbox{-}}
\setlength{\emergencystretch}{2em}
\multlinegap0pt
\usepackage{amssymb}
\usepackage{bbm}
\usepackage{enumerate}
\usepackage[shortlabels]{enumitem}
\usepackage{tikz}
\newtheorem{lemma}{Lemma}
\newcommand{\bs}{\boldsymbol}
\title{Exact analysis of transient behavior\\ of finite-capacity MAP-driven queues}
\author{Michel Mandjes, Dani\"el Rutgers, Werner Scheinhardt}
\date{Source: arXiv:2602.09676v1 (CC BY 4.0)}
\begin{document}
\maketitle

\noindent\emph{Source and licence.} Exact analysis of transient behavior\\ of finite-capacity MAP-driven queues. Version arXiv:2602.09676v1 (CC BY 4.0), distributed by arXiv under the Creative Commons Attribution 4.0 International licence, http://creativecommons.org/licenses/by/4.0/, as recorded in the arXiv licence metadata for this exact version and shown on the abstract page https://arxiv.org/abs/2602.09676v1. Full licence notice: \texttt{licenses/arxiv-2602.09676-CC-BY-4.0.txt}.\par\smallskip

\noindent\textit{Editorial scope.} Counted unit: the article's lemma delta* (source0.tex lines 638--642). The article's complete original proof of it is delivered with it (source0.tex lines 643--685, one \texttt{proof} environment of 43 lines). No mathematics is written, repaired or re-derived: the statement and the proof are the article's own lines, in the article's own notation, and no retained line is substituted or re-flowed.  No further source-local context is needed: every cross-reference of every retained line resolves inside the two delivered spans.  Nothing was omitted from any retained span, so the package carries no omission record.  No retained line contains a citation, so the wrapper carries no bibliography and no bibliography entry.  No retained line contains a figure environment, an image-inclusion command or drawing code, so no image is delivered and the package declares no asset dependency.  Editorial formatting only, in addition: an AMS \texttt{amsart} preamble that restates the article's own declarations and macros used by the retained lines, namely the environments \texttt{lemma} on separate counters, together with the standard packages those lines need.  The retained environments print in this document as Lemma 1 in order of appearance, whereas the article prints the same items with the numbering of its own full document.  The complete native source of the article as distributed in the arXiv e-print for this exact version is bundled unchanged as \texttt{sources/194336/source0.tex}.
\begin{lemma}\label{lemma: delta*} For $\alpha \geqslant0, \beta>0$ and $x\in[0,K]$,
    \begin{equation}\label{eq: delta*}
        \bs\delta_\star(x) = e^{-\alpha x}\left( \bs{I} - e^{\alpha x}\bs\delta_-(x) - e^{-\alpha(K-x)}\bs\eta(K-x,\alpha,\beta) + e^{-\alpha(K-x)}\bs\delta_-(x)\,\bs\eta(K,\alpha,\beta) \right)\bs\Phi(\alpha,\beta).
    \end{equation}
\end{lemma}

\begin{proof} 
    We start by decomposing $\bs\delta_\star(x)$. It is readily verified that
   \begin{equation}\label{eq: decomp delta*}
       \bs\delta_\star(x) = e^{-\alpha x}\left(\bs\Phi(\alpha,\beta) - {\bs \Xi}(x,\alpha, \beta) -{\bs \Psi}(x,\alpha, \beta)\right),
   \end{equation}
    where the matrices ${\bs \Xi}(x,\alpha,\beta)$ and ${\bs \Psi}(x,\alpha,\beta)$ have entries
\begin{align*}
    \Xi_{ij}(x,\alpha,\beta)&:= {\mathbb E}_i\left(e^{-\alpha Y(T_{\beta})}{\bs 1}_{\{\sigma(x)\leqslant\min\{\tau(K-x), T_{\beta} \},\,J(T_{\beta})=j\}}\right),\\
    \Psi_{ij}(x,\alpha,\beta)&:= {\mathbb E}_i\left(e^{-\alpha Y(T_{\beta})}{\bs 1}_{\{\tau(K-x)\leqslant\min\{\sigma(x), T_{\beta} \},\,J(T_{\beta})=j\}}\right).
\end{align*}Hence, to analyze $\bs\delta_\star(x)$, we may consider all sample paths of the free process $Y$ except those on which the level $-x$ or $K-x$ is reached first, before $T_\beta$. We proceed by analyzing (i)~${\bs \Xi}(x,\alpha,\beta)$ and (ii)~${\bs \Psi}(x,\alpha,\beta)$.

\medskip
(i)~By the strong Markov property, we find that the entries of ${\bs \Xi}(x,\alpha,\beta)$ are simply given by
\begin{align*}
    \Xi_{ij}(x,\alpha,\beta) = e^{\alpha x}\sum_{k=1}^d \delta_{-,ik}(x)\,{\Phi}_{kj}(\alpha, \beta).
\end{align*}
The reasoning for this decomposition is that in this scenario the level $-x$ is reached (with equality, due to the absence of negative jumps) before $T_{\beta}$ and $\tau(K-x)$, where the summation over $k$ takes care of the possible states of the background process at time $\sigma(x)$.

\medskip
(ii)~We further decompose ${\bs \Psi}(x,\alpha,\beta)$ by considering all paths in which the level $K-x$ is reached before $T_{\beta}$ except those in which the level $-x$ is reached before the level $K-x$. This gives
\begin{equation}\label{eq: decomp Psi}
    {\Psi}_{ij}(x,\alpha,\beta) = {\mathbb E}_i\left(e^{-\alpha Y(T_{\beta})}{\bs 1}_{\{\tau(K-x)\leqslant T_{\beta},\,J(T_{\beta})=j\}}\right) - {\mathbb E}_i\left(e^{-\alpha Y(T_{\beta})}{\bs 1}_{\{\sigma(x)\leqslant\tau(K-x)\leqslant T_{\beta},\,J(T_{\beta})=j\}}\right).
\end{equation}
The first component in the right-hand side of \eqref{eq: decomp Psi} can be rewritten as 
\begin{equation*}
    {\mathbb E}_i\left(e^{-\alpha Y(T_{\beta})}{\bs 1}_{\{\tau(K-x)\leqslant T_{\beta},\,J(T_{\beta})=j\}}\right) = e^{-\alpha(K-x)}\sum_{k=1}^d  \eta_{ik}(K-x, \alpha, \beta)\,\Phi_{kj}(\alpha,\beta),
\end{equation*}
since in this case the level $K-x$ is reached before $T_{\beta}$, with $k$ and $\eta_{ik}(K-x,\alpha,\beta)$ accounting for the background state and the overshoot at time $\tau(K-x)$, respectively.

The second component in the right-hand side of \eqref{eq: decomp Psi} can be rewritten as
\begin{equation*}
        {\mathbb E}_i\left(e^{-\alpha Y(T_{\beta})}{\bs 1}_{\{\sigma(x)\leqslant\tau(K-x)\leqslant T_{\beta},\,J(T_{\beta})=j\}}\right) = \sum_{k=1}^d\sum_{\ell=1}^d\delta_{-,ik}(x)\,e^{-\alpha(K-x)}\eta_{k\ell}(K,\alpha,\beta)\,\Phi_{\ell j}(\alpha,\beta),
    \end{equation*}
where we argue that, in this case, the paths first reach level $-x$ (with equality) and subsequently reach level $K-x$ (or, equivalently, level $K$ in a system shifted by $x$ and starting at $0$) before $T_{\beta}$.

\medskip

Inserting the found expressions for ${\bs \Xi}(x,\alpha,\beta)$ and ${\bs \Psi}(x,\alpha,\beta)$ into \eqref{eq: decomp delta*}, we obtain the following matrix equation for $\bs{\Psi}(x,\alpha,\beta)$:
\begin{equation*}
    {\bs \Psi}(x,\alpha,\beta) = e^{-\alpha(K-x)}\left(\bs{\eta}(K-x, \alpha, \beta) - \bs{\delta}_{-}(x)\,\bs{\eta}(K,\alpha,\beta)\right)\,\bs{\Phi}(\alpha,\beta)
\end{equation*}
Applying the identity \eqref{eq: decomp delta*}, we can now determine $\bs\delta_{\star}(x)$ in a straightforward manner.
\end{proof}

\end{document}
